\ifdefined\gkver\else\def\gkver{student}\fi
\documentclass[\gkver]{gaokaozhenti}
\begin{document}

\chapter{2009年江苏}

\begin{enumerate}
\item
%% number: 1
%% paperName: 2009年江苏高考第 1 题 3 分
%% typeId: 1
%% type: 单选题
%% chapter: 电场
%% point: 库仑定律
%% method: 无
%% score: 3
%% degree: 650
%% duplicateId: 0
%% body:
两个分别带有电荷量 $-Q$ 和 $+3Q$ 的相同金属小球（均可视为点电荷），固定在相距为 $r$ 的两处，它们间库仑力的大小为 $F$。两小球相互接触后将其固定距离变为 $r/2$，则两球间库仑力的大小为\xzanswer{C}
\fourchoices[answer=C]
{$\frac{1}{12}F$}
{$\frac{3}{4}F$}
{$\frac{4}{3}F$}
{$12F$}
%% answer: C
%% memo:
\memoanswer{
	本题考查库仑定律及带电体电量的转移问题。接触前两个点电荷之间的库仑力大小为 $F=k\frac{Q\cdot 3Q}{r^{2}}$，两个相同的金属球各自带电，接触后再分开，其所带电量先中和后均分，所以两球分开后各自带电为 $+Q$，距离又变为原来的 $\frac{1}{2}$，库仑力为 $F'=k\frac{Q\cdot Q}{\left(\frac{r}{2}\right)^{2}}$，所以两球间库仑力的大小为 $\frac{4}{3}F$，C 项正确。如两球原来带正电，则接触各自带电均为 $+2Q$。
}

\item
%% number: 2
%% paperName: 2009年江苏高考第 2 题 3 分
%% typeId: 1
%% type: 单选题
%% chapter: 力物体的平衡
%% point: 三力平衡
%% method: 无
%% score: 3
%% degree: 650
%% duplicateId: 0
%% body:
用一根长 $1\Um$ 的轻质细绳将一幅质量为 $1\Ukg$ 的画框对称悬挂在墙壁上，已知绳能承受的最大张力为 $10\UN$，为使绳不断裂，画框上两个挂钉的间距最大为（$g$ 取 $10\Umsq$）\xzanswer{A}
\onepicture[w=3.6cm]{figs/02.png}
\fourchoices[answer=A]
{$\frac{\sqrt{3}}{2}\Um$}
{$\frac{\sqrt{2}}{2}\Um$}
{$\frac{1}{2}\Um$}
{$\frac{\sqrt{3}}{4}\Um$}
%% answer: A
%% memo:
\memoanswer{
	熟练应用力的合成和分解以及合成与分解中的一些规律，是解决本题的根本；一个大小方向确定的力分解为两个等大的力时，合力在分力的角平分线上，且两分力的夹角越大，分力越大。题中当绳子拉力达到 $F=10\UN$ 的时候，绳子间的张角最大，即两个挂钉间的距离最大；画框受到重力和绳子的拉力，三个力为共点力。绳子与竖直方向的夹角为 $\theta$，绳子长为 $L_{0}=1\Um$，则有 $mg=2F\cos\theta$，两个挂钉间的距离 $L=2\cdot\frac{L_{0}}{2}\sin\theta$，解得 $L=\frac{\sqrt{3}}{2}\Um$，A 项正确。
}

\item
%% number: 3
%% paperName: 2009年江苏高考第 3 题 3 分
%% typeId: 1
%% type: 单选题
%% chapter: 曲线运动
%% point: 天体运动
%% method: 无
%% score: 3
%% degree: 450
%% duplicateId: 0
%% body:
英国《新科学家（New Scientist）》杂志评选出了 2008 年度世界 8 项科学之最，在 XTEJ1650-500 双星系统中发现的最小黑洞位列其中，若某黑洞的半径 $R$ 约 $45\Ukm$，质量 $M$ 和半径 $R$ 的关系满足 $\frac{M}{R}=\frac{c^{2}}{2G}$（其中 $c$ 为光速，$G$ 为引力常量），则该黑洞表面重力加速度的数量级为\xzanswer{C}
\fourchoices[answer=C]
{$10^{8}\Umsq$}
{$10^{10}\Umsq$}
{$10^{12}\Umsq$}
{$10^{14}\Umsq$}
%% answer: C
%% memo:
\memoanswer{
	处理本题要从所给的材料中，提炼出有用信息，构建好物理模型，选择合适的物理方法求解。黑洞实际为一天体，天体表面的物体受到的重力近似等于物体与该天体之间的万有引力，对黑洞表面的某一质量为 $m$ 的物体有 $G\frac{Mm}{R^{2}}=mg$，又有 $\frac{M}{R}=\frac{c^{2}}{2G}$，联立解得 $g=\frac{c^{2}}{2R}$，代入数据得重力加速度的数量级为 $10^{12}\Umsq$，C 项正确。
}

\item
%% number: 4
%% paperName: 2009年江苏高考第 4 题 3 分
%% typeId: 1
%% type: 单选题
%% chapter: 曲线运动
%% point: 平抛运动
%% method: 无
%% score: 3
%% degree: 450
%% duplicateId: 0
%% body:
在无风的情况下，跳伞运动员从水平飞行的飞机上跳伞，下落过程中受到空气阻力，下列描绘下落速度的水平分量大小 $v_{x}$、竖直分量大小 $v_{y}$ 与时间 $t$ 的图像，可能正确的是\xzanswer{B}
\fourchoices[answer=B, ispicture=true, h=2.2cm]
{figs/04a.png}
{figs/04b.png}
{figs/04c.png}
{figs/04d.png}
%% answer: B
%% memo:
\memoanswer{
	跳伞运动员下落过程中受到的空气阻力并非为恒力，与速度有关，且速度越大受到的阻力越大，知道速度与所受阻力的规律是解决本题的关键。竖直方向运动员受重力和空气阻力，速度逐渐增大，阻力增大合力减小，加速度减小；水平方向只受阻力，速度减小，阻力减小，加速度减小。在 $v-t$ 图象中图线的斜率表示加速度，B 项正确。
}

\item
%% number: 5
%% paperName: 2009年江苏高考第 5 题 3 分
%% typeId: 1
%% type: 单选题
%% chapter: 电路
%% point: 电容
%% method: 无
%% score: 3
%% degree: 450
%% duplicateId: 0
%% body:
在如图所示的闪光灯电路中，电源的电动势为 $E$，电容器的电容为 $C$。当闪光灯两端电压达到击穿电压 $U$ 时，闪光灯才有电流通过并发光，正常工作时，闪光灯周期性短暂闪光，则可以判定\xzanswer{D}
\onepicture[w=5.0cm]{figs/05.png}
\fourchoices[answer=D]
{电源的电动势 $E$ 一定小于击穿电压 $U$}
{电容器所带的最大电荷量一定为 $CE$}
{闪光灯闪光时，电容器所带的电荷量一定增大}
{在一个闪光周期内，通过电阻 $R$ 的电荷量与通过闪光灯的电荷量一定相等}
%% answer: D
%% memo:
\memoanswer{
	理解此电路的工作过程是解决本题的关键。电容器两端的电压与闪光灯两端的电压相等，当电源给电容器充电，达到闪光灯击穿电压 $U$ 时，闪光灯被击穿，电容器放电，放电后闪光灯两端电压小于 $U$，断路，电源再次给电容器充电，达到电压 $U$ 时，闪光灯又被击穿，电容器放电，如此周期性充放电，使得闪光灯周期性短暂闪光。要使得充电后达到电压 $U$，则电源电动势一定大于等于 $U$，A 项错误；电容器两端的最大电压为 $U$，故电容器所带的最大电荷量为 $CU$，B 项错误；闪光灯闪光时电容器放电，所带电荷量减少，C 项错误；充电时电荷通过 $R$，通过闪光灯放电，故充放电过程中通过电阻 $R$ 的电荷量与通过闪光灯的电荷量一定相等，D 项正确。
}

\item
%% number: 6
%% paperName: 2009年江苏高考第 6 题 4 分
%% typeId: 2
%% type: 多选题
%% chapter: 交流电
%% point: 变压器
%% method: 无
%% score: 4
%% degree: 450
%% duplicateId: 0
%% body:
如图所示，理想变压器的原、副线圈匝数比为 1:5，原线圈两端的交变电压为 $u=20\sqrt{2}\sin100\pi t\UV$，氖泡在两端电压达到 $100\UV$ 时开始发光，下列说法中正确的有\xzanswer{AB}
\onepicture[w=5.6cm]{figs/06.png}
\fourchoices[answer=AB]
{开关接通后，氖泡的发光频率为 $100\UHz$}
{开关接通后，电压表的示数为 $100\UV$}
{开关断开后，电压表的示数变大}
{开关断开后，变压器的输出功率不变}
%% answer: AB
%% memo:
\memoanswer{
	本题主要考查变压器的知识，要能对变压器的最大值、有效值、瞬时值以及变压器变压原理、功率等问题彻底理解。由交变电压的瞬时值表达式知，原线圈两端电压的有效值为 $U_{1}=\frac{20\sqrt{2}}{\sqrt{2}}\UV=20\UV$，由 $\frac{n_{1}}{n_{2}}=\frac{U_{1}}{U_{2}}$ 得副线圈两端的电压为 $U_{2}=100\UV$，电压表的示数为交流电的有效值，B 项正确；交变电压的频率为 $f=\frac{100\pi}{2\pi}\UHz=50\UHz$，一个周期内电压两次大于 $100\UV$，即一个周期内氖泡能两次发光，所以其发光频率为 $100\UHz$，A 项正确；开关断开前后，输入电压不变，变压器的变压比不变，故输出电压不变，C 项错误；断开后，电路消耗的功率减小，输出功率决定输入功率，D 项错误。
}

\item
%% number: 7
%% paperName: 2009年江苏高考第 7 题 4 分
%% typeId: 2
%% type: 多选题
%% chapter: 直线运动
%% point: 匀变速直线运动
%% method: 无
%% score: 4
%% degree: 450
%% duplicateId: 0
%% body:
如图所示，以 $8\Ums$ 匀速行驶的汽车即将通过路口，绿灯还有 $2\Us$ 将熄灭，此时汽车距离停车线 $18\Um$。该车加速时最大加速度大小为 $2\Umsq$，减速时最大加速度大小为 $5\Umsq$。此路段允许行驶的最大速度为 $12.5\Ums$，下列说法中正确的有\xzanswer{AC}
\onepicture[w=5.0cm]{figs/07.png}
\fourchoices[answer=AC]
{如果立即做匀加速运动，在绿灯熄灭前汽车可能通过停车线}
{如果立即做匀加速运动，在绿灯熄灭前通过停车线汽车一定超速}
{如果立即做匀减速运动，在绿灯熄灭前汽车一定不能通过停车线}
{如果距停车线 $5\Um$ 处减速，汽车能停在停车线处}
%% answer: AC
%% memo:
\memoanswer{
	熟练应用匀变速直线运动的公式，是处理问题的关键，对汽车运动的问题一定要注意所求解的问题是否与实际情况相符。如果立即做匀加速直线运动，$t_{1}=2\Us$ 内的位移 $x=v_{0}t_{1}+\frac{1}{2}a_{1}t_{1}^{2}=20\Um>18\Um$，此时汽车的速度为 $v_{1}=v_{0}+a_{1}t_{1}=12\Ums<12.5\Ums$，汽车没有超速，A 项正确；如果立即做匀减速运动，速度减为零需要时间 $t_{2}=\frac{v_{0}}{a_{2}}=1.6\Us$，此过程通过的位移为 $x_{2}=\frac{1}{2}a_{2}t_{2}^{2}=6.4\Um$，C 项正确、D 项错误。
}

\item
%% number: 8
%% paperName: 2009年江苏高考第 8 题 4 分
%% typeId: 2
%% type: 多选题
%% chapter: 电场
%% point: 电场中的图象
%% method: 无
%% score: 4
%% degree: 450
%% duplicateId: 0
%% body:
空间某一静电场的电势 $\varphi$ 在 $x$ 轴上分布如图所示，$x$ 轴上两点 B、C 点电场强度在 $x$ 方向上的分量分别是 $E_{Bx}$、$E_{Cx}$，下列说法中正确的有\xzanswer{AD}
\onepicture[w=5.4cm]{figs/08.png}
\fourchoices[answer=AD]
{$E_{Bx}$ 的大小大于 $E_{Cx}$ 的大小}
{$E_{Bx}$ 的方向沿 $x$ 轴正方向}
{电荷在 O 点受到的电场力在 $x$ 方向上的分量最大}
{负电荷沿 $x$ 轴从 B 移到 C 的过程中，电场力先做正功，后做负功}
%% answer: AD
%% memo:
\memoanswer{
	本题的入手点在于如何判断 $E_{Bx}$ 和 $E_{Cx}$ 的大小，由图象可知在 $x$ 轴上各点的电场强度在 $x$ 方向的分量不相同，如果在 $x$ 方向上取极小的一段，可以把此段看做是匀强电场，用匀强电场的处理方法思考，从而得到结论，此方法为微元法；需要对电场力的性质和能的性质有较为全面的理解。在 B 点和 C 点附近分别取很小的一段 $d$，由图象，B 点段对应的电势差大于 C 点段对应的电势差，看做匀强电场有 $E=\frac{\Delta\varphi}{d}$，可见 $E_{Bx}>E_{Cx}$，A 项正确；同理可知 O 点场强最小，电荷在该点受到的电场力最小，C 项错误；沿电场方向电势降低，在 O 点左侧，$E_{Bx}$ 的方向沿 $x$ 轴负方向，在 O 点右侧，$E_{Cx}$ 的方向沿 $x$ 轴正方向，所以 B 项错误，D 项正确。
}

\item
%% number: 9
%% paperName: 2009年江苏高考第 9 题 4 分
%% typeId: 2
%% type: 多选题
%% chapter: 运动和力
%% point: 变加速直线运动
%% method: 无
%% score: 4
%% degree: 250
%% duplicateId: 0
%% body:
如图所示，两质量相等的物块 A、B 通过一轻质弹簧连接，B 足够长、放置在水平面上，所有接触面均光滑。弹簧开始时处于原长，运动过程中始终处在弹性限度内。在物块 A 上施加一个水平恒力，A、B 从静止开始运动到第一次速度相等的过程中，下列说法中正确的有\xzanswer{BCD}
\onepicture[w=5.6cm]{figs/09.png}
\fourchoices[answer=BCD]
{当 A、B 加速度相等时，系统的机械能最大}
{当 A、B 加速度相等时，A、B 的速度差最大}
{当 A、B 的速度相等时，A 的速度达到最大}
{当 A、B 的速度相等时，弹簧的弹性势能最大}
%% answer: BCD
%% memo:
\memoanswer{
	处理本题的关键是对物体进行受力分析和运动过程分析，使用图象处理则可以使问题大大简化。对 A、B 在水平方向受力分析如图，$F_{1}$ 为弹簧的拉力；当加速度大小相同为 $a$ 时，对 A 有 $F-F_{1}=ma$，对 B 有 $F_{1}=ma$，得 $F_{1}=\frac{F}{2}$，在整个过程中 A 的合力（加速度）一直减小而 B 的合力（加速度）一直增大，在达到共同加速度之前 A 的合力（加速度）一直大于 B 的合力（加速度），之后 A 的合力（加速度）一直小于 B 的合力（加速度）。
	\threepicture[showlabel=false, wA=4.2cm, wB=3.2cm, wC=6.0cm]
	{figs/09_an1.jpg}
	{figs/09_an2.jpg}
	{figs/09_an3.jpg}
	两物体运动的 $v-t$ 图象如图，$t_{1}$ 时刻，两物体加速度相等，斜率相同，速度差最大；$t_{2}$ 时刻两物体的速度相等，A 速度达到最大值，两实线之间围成的面积有最大值即两物体的相对位移最大，弹簧被拉到最长；除重力和弹簧弹力外其它力对系统做正功，系统机械能增加，$t_{1}$ 时刻之后拉力依然做正功，即加速度相等时，系统机械能并非最大值。
}

\item
%% number: 10
%% paperName: 2009年江苏高考第 10 题 8 分
%% typeId: 1
%% type: 单选题
%% chapter: 电路
%% point: 伏安法测电阻
%% method: 无
%% score: 8
%% degree: 850
%% duplicateId: 0
%% body:
有一根圆台状均匀质合金棒如图 \ref{2009江苏10a} 所示，某同学猜测其电阻的大小与该合金棒的电阻率 $\rho$、长度 $L$ 和两底面直径 $d$、$D$ 有关。他进行了如下实验：

（1）用游标卡尺测量合金棒的两底面直径 $d$、$D$ 和长度 $L$。图 \ref{2009江苏10b} 中游标卡尺（游标尺上有 20 个等分刻度）的读数 $L=$\tkanswer[2.2cm]{9.940} cm。

（2）测量该合金棒电阻的实物电路如图 \ref{2009江苏10c} 所示（相关器材的参数已在图中标出）。该合金棒的电阻约为几个欧姆。图中有一处连接不当的导线是\tkanswer[1.2cm]{⑥}。（用标注在导线旁的数字表示）

（3）改正电路后，通过实验测得合金棒的电阻 $R=6.72\UO$。根据电阻定律计算电阻率为 $\rho$、长为 $L$、直径分别为 $d$ 和 $D$ 的圆柱状合金棒的电阻分别为 $R_{d}=13.3\UO$、$R_{D}=3.38\UO$。他发现：在误差允许范围内，电阻 $R$ 满足 $R^{2}=R_{d}R_{D}$，由此推断该圆台状合金棒的电阻 $R=$\tkanswer[3.2cm]{$\frac{4\rho L}{\pi dD}$}。（用 $\rho$、$L$、$d$、$D$ 表述）
\threepicture[labelA=2009江苏10a, labelB=2009江苏10b, labelC=2009江苏10c, numstyle=chinese, wA=3.4cm, wB=4.6cm, wC=5.6cm, align=right]
{figs/10a.png}
{figs/10b.png}
{figs/10c.png}
%% answer: D
%% memo:
\memoanswer{
	（1）游标卡尺的读数，按步骤进行则不会出错。首先，确定游标卡尺的精度为 20 分度，即为 $0.05\Umm$，然后以毫米为单位从主尺上读出整毫米数 $99.00\Umm$，注意小数点后的有效数字要与精度一样，再从游标尺上找出对的最齐的一根刻线，精度×格数 $=0.05\times8\Umm=0.40\Umm$，最后两者相加，根据题目单位要求换算为需要的数据，$99.00\Umm+0.40\Umm=99.40\Umm=9.940\Ucm$。
	（2）本实验为测定一个几欧姆的电阻，在用伏安法测量其两端的电压和通过电阻的电流时，因为安培表的内阻较小，为了减小误差，应用安培表外接法，⑥线的连接使用的是安培表内接法。
	（3）审题是处理本题的关键，弄清题意也就能够找到处理本题的方法。根据电阻定律计算电阻率为 $\rho$、长为 $L$、直径分别为 $d$ 和 $D$ 的圆柱状合金棒的电阻分别为 $R_{d}=13.3\UO$、$R_{D}=3.38\UO$，即 $R_{d}=\rho\frac{L}{\pi\left(\frac{d}{2}\right)^{2}}$，$R_{D}=\rho\frac{L}{\pi\left(\frac{D}{2}\right)^{2}}$，而电阻 $R$ 满足 $R^{2}=R_{d}R_{D}$，将 $R_{d}$、$R_{D}$ 代入得 $R=\frac{4\rho L}{\pi dD}$。
}

\item
%% number: 11
%% paperName: 2009年江苏高考第 11 题 10 分
%% typeId: 4
%% type: 实验题
%% chapter: 运动和力
%% point: 验证牛顿第二定律
%% method: 无
%% score: 10
%% degree: 450
%% duplicateId: 0
%% body:
“探究加速度与物体质量、物体受力的关系”的实验装置如图 \ref{2009江苏11a} 所示。

\begin{enumerate}
	\item
	在平衡小车与桌面之间摩擦力的过程中，打出了一条纸带如图 \ref{2009江苏11b} 所示。计时器打点的时间间隔为 $0.02\Us$。从比较清晰的点起，每 5 个点取一个计数点，量出相邻计数点之间的距离。该小车的加速度 $a=$\tkanswer[2.0cm]{} $\Umsq$。（结果保留两位有效数字）
	\item
	平衡摩擦力后，将 5 个相同的砝码都放在小车上。挂上砝码盘，然后每次从小车上取一个砝码添加到砝码盘中，测量小车的加速度。小车的加速度 $a$ 与砝码盘中砝码总重力 $F$ 的实验数据如下表：

	\begin{center}
	\begin{tabular}{|c|c|c|c|c|c|}
	\hline
	砝码盘中砝码总重力 $F/\UN$ & 0.196 & 0.392 & 0.588 & 0.784 & 0.980 \\
	\hline
	加速度 $a/(\Um\cdot\Us^{-2})$ & 0.69 & 1.18 & 1.66 & 2.18 & 2.70 \\
	\hline
	\end{tabular}
	\end{center}

	请根据实验数据作出 $a-F$ 的关系图像。

	\onepicture[w=5.0cm, align=right]{figs/11c.png}
	\item
	根据提供的实验数据作出的 $a-F$ 图线不通过原点，请说明主要原因。
\end{enumerate}
\twopicture[labelA=2009江苏11a, labelB=2009江苏11b, numstyle=chinese, wA=5.0cm, wB=6.0cm, align=right]
{figs/11a.png}
{figs/11b.png}
%% answer:
\jdanswer{
\begin{enumerate}
	\item
	$a=0.16\Umsq$（或 $0.15\Umsq$）
	\item
	图见解析。
	\item
	未计入砝码盘的重力。
\end{enumerate}
}
%% memo:
\memoanswer{
	（1）处理匀变速直线运动中所打出的纸带，求解加速度用公式 $\Delta x=aT^{2}$，关键弄清公式中各个量的物理意义，$\Delta x$ 为连续相等时间内的位移差，$T$ 为连续相等的时间间隔，如果每 5 个点取一个点，则连续两点的时间间隔为 $T=0.1\Us$，$\Delta x=(3.68-3.52)\times10^{-2}\Um$，代入可得加速度 $a=0.16\Umsq$。也可以使用最后一段和第二段的位移差求解，得加速度 $a=0.15\Umsq$。
	（2）根据图中的数据，合理的设计横纵坐标的刻度值，使图线倾斜程度不能太小也不能太大，以与水平方向夹角 $45^{\circ}$ 左右为宜。由此确定 $F$ 的范围从 0 设置到 $1\UN$ 较合适，而 $a$ 则从 0 到 $3\Umsq$ 较合适。设好刻度，根据数据确定各点的位置，将各点用一条直线连起来，延长交与坐标轴某一点。如图所示。
	\onepicture[w=6.0cm, align=right]{figs/11_an.jpg}
	（3）处理图象问题要注意图线的斜率、交点、拐点、面积等意义，能正确理解这些量的意义则很多问题将会迎刃而解。与纵坐标相交而不过原点，该交点说明当不挂砝码时，小车仍有加速度，即绳对小车仍有拉力，从此拉力的来源考虑很容易得到答案，是因为砝码盘的重力，而在（2）问的图表中只给出了砝码的总重力，而没有考虑砝码盘的重力。
}

\item
%% number: 12
%% paperName: 2009年江苏高考第 12 题 12 分
%% typeId: 3
%% type: 填空题
%% chapter: 原子和原子核
%% point: 人工核反应
%% method: 无
%% score: 12
%% degree: 450
%% duplicateId: 0
%% body:
本题为选做题，包括 A、B、C 三小题，请选定其中两题，并在相应的答题区域内作答。若三题都做，则按 A、B 两题评分。

\textbf{A．}（选修模块 3-3）（12 分）

\begin{enumerate}
	\item
	若一气泡从湖底上升到湖面的过程中温度保持不变，则在此过程中关于气泡中的气体，下列说法正确的是\xzanswer{D}（填写选项前的字母）
	\fourchoices[answer=D]
	{气体分子间的作用力增大}
	{气体分子的平均速率增大}
	{气体分子的平均动能减小}
	{气体组成的系统的熵增加}
	\item
	若将气泡内的气体视为理想气体，气泡从湖底上升到湖面的过程中，对外界做了 $0.6\UJ$ 的功，则此过程中的气泡\tkanswer[1.2cm]{吸收}（填“吸收”或“放出”）的热量是\tkanswer[1.4cm]{$0.6\UJ$}。气泡到达湖面后，温度上升的过程中，又对外界做了 $0.1\UJ$ 的功，同时吸收了 $0.3\UJ$ 的热量，则此过程中，气泡内气体内能增加了\tkanswer[1.4cm]{$0.2\UJ$}。
	\item
	已知气泡内气体的密度为 $1.29\Ukgmc$，平均摩尔质量为 $0.29\Ukgmol$。阿伏加德罗常数 $N_{A}=6.02\times10^{23}\Umol^{-1}$，取气体分子的平均直径为 $2\times10^{-10}\Um$，若气泡内的气体能完全变为液体，请估算液体体积与原来气体体积的比值。（结果保留一位有效数字）
\end{enumerate}

\textbf{B．}（选修模块 3-4）（12 分）

\begin{enumerate}
	\item
	如图所示，强强乘速度为 $0.9c$（$c$ 为光速）的宇宙飞船追赶正前方的壮壮，壮壮的飞行速度为 $0.5c$，强强向壮壮发出一束光进行联络，则壮壮观测到该光束的传播速度为\xzanswer{D}（填写选项前的字母）
	\onepicture[w=5.4cm]{figs/12a.png}
	\fourchoices[answer=D]
	{$0.4c$}
	{$0.5c$}
	{$0.9c$}
	{$1.0c$}
	\item
	在 $t=0$ 时刻，质点 A 开始做简谐运动，其振动图象如图所示。质点 A 振动的周期是\tkanswer[1.0cm]{$4\Us$}；$t=8\Us$ 时，质点 A 的运动沿 $y$ 轴的\tkanswer[1.0cm]{正}方向（填“正”或“负”）；质点 B 在波动的传播方向上与 A 相距 $16\Um$，已知波的传播速度为 $2\Ums$，在 $t=9\Us$ 时，质点 B 偏离平衡位置的位移是\tkanswer[1.2cm]{$10\Ucm$}。
	\onepicture[w=5.6cm]{figs/12b.png}
	\item
	如图所示是北京奥运会期间安置在游泳池底部的照相机拍摄的一张照片，照相机的镜头竖直向上。照片中，水立方运动馆的景象呈现在半径 $r=11\Ucm$ 的圆形范围内，水面上的运动员手到脚的长度 $l=10\Ucm$，若已知水的折射率为 $n=\frac{4}{3}$，请根据运动员的实际身高估算该游泳池的水深 $h$。（结果保留两位有效数字）
	\onepicture[w=4.6cm]{figs/12c.png}
\end{enumerate}

\textbf{C．}（选修模块 3-5）（12 分）

在 $\beta$ 衰变中常伴有一种称为“中微子”的粒子放出。中微子的性质十分特别，因此在实验中很难探测。1953 年，莱尼斯和柯文建造了一个由大水槽和探测器组成的实验系统，利用中微子与水中 $\ce{^{1}_{1}H}$ 的核反应，间接地证实了中微子的存在。

\begin{enumerate}
	\item
	中微子与水中的 $\ce{^{1}_{1}H}$ 发生核反应，产生中子（$\ce{^{1}_{0}n}$）和正电子（$\ce{^{0}_{1}e}$），即 中微子 $+\ce{^{1}_{1}H -> ^{1}_{0}n + ^{0}_{1}e}$ 可以判定，中微子的质量数和电荷数分别是\xzanswer{A}（填写选项前的字母）
	\fourchoices[answer=A]
	{0 和 0}
	{0 和 1}
	{1 和 0}
	{1 和 1}
	\item
	上述核反应产生的正电子与水中的电子相遇，与电子形成几乎静止的整体后，可以转变为两个光子（$\gamma$），即 $\ce{^{0}_{1}e + ^{0}_{-1}e -> 2\gamma}$。已知正电子和电子的质量都为 $9.1\times10^{-31}\Ukg$，反应中产生的每个光子的能量约为\tkanswer[1.8cm]{$8.2\times10^{-14}\UJ$}。正电子与电子相遇不可能只转变为一个光子，原因是\tkanswer[4.0cm]{此过程遵循动量守恒}。
	\item
	试通过分析比较，具有相同动能的中子和电子的物质波波长的大小。
\end{enumerate}
%% answer:
\jdanswer{
\begin{enumerate}
	\item
	A：（1）D；（2）吸收，$0.6\UJ$，$0.2\UJ$；（3）$1\times10^{-4}$（$9\times10^{-5}\sim2\times10^{-4}$ 都算对）。
	\item
	B：（1）D；（2）$4\Us$，正，$10\Ucm$；（3）$2.1\Um$（$1.6\sim2.6\Um$ 都算对）。
	\item
	C：（1）A；（2）$8.2\times10^{-14}\UJ$，遵循动量守恒；（3）$\lambda_{n}<\lambda_{e}$。
\end{enumerate}
}
%% memo:
\memoanswer{
	\textbf{A．}（1）掌握分子动理论和热力学定律才能准确处理本题。气泡的上升过程气泡内的压强减小，温度不变，由玻意耳定律知，上升过程中体积增大，微观上体现为分子间距增大，分子间引力减小，温度不变所以气体分子的平均动能、平均速率不变，此过程为自发过程，故熵增大。D 项正确。

	（2）本题从热力学第一定律入手，抓住理想气体内能只与温度有关的特点进行处理。理想气体等温过程中内能不变，由热力学第一定律 $\Delta U=Q+W$，物体对外做功 $0.6\UJ$，则一定同时从外界吸收热量 $0.6\UJ$，才能保证内能不变。而温度上升的过程，内能增加了 $0.2\UJ$。

	（3）设气体体积为 $V_{0}$，液体体积为 $V_{1}$。气体分子数 $n=\frac{\rho V_{0}}{m}N_{A}$，$V_{1}=n\frac{\pi d^{3}}{6}$（或 $V_{1}=nd^{3}$），则 $\frac{V_{1}}{V_{0}}=\frac{\rho}{6m}\pi d^{3}N_{A}$（或 $\frac{V_{1}}{V_{0}}=\frac{\rho}{m}d^{3}N_{A}$），解得 $\frac{V_{1}}{V_{0}}=1\times10^{-4}$（$9\times10^{-5}\sim2\times10^{-4}$ 都算对）。

	\textbf{B．}（1）根据爱因斯坦相对论，在任何参考系中，光速不变。D 项正确。

	（2）振动图象和波形图比较容易混淆，而导致出错，在读图时一定要注意横纵坐标的物理意义，以避免出错。题图为波的振动图象，由图象可知周期为 $4\Us$，波源的起振方向与波头的振动方向相同且向上，$t=6\Us$ 时质点在平衡位置向下振动，故 $8\Us$ 时质点在平衡位置向上振动；波传播到 B 点，需要时间 $t_{1}=\frac{x}{v}=\frac{16}{2}\Us=8\Us$，故 $t=9\Us$ 时，质点又振动了 $1\Us$（$\frac{1}{4}$ 个周期），处于正向最大位移处，位移为 $10\Ucm$。

	（3）根据题意能画出光路图，正确使用物像比是解决本题的关键。设照片圆形区域的实际半径为 $R$，运动员的实际长为 $L$，光路如图所示。折射定律 $n\sin\alpha=\sin90^{\circ}$，几何关系 $\sin\alpha=\frac{R}{\sqrt{R^{2}+h^{2}}}$，物像比 $\frac{R}{r}=\frac{L}{l}$，得 $h=\sqrt{n^{2}-1}\cdot\frac{L}{l}r$，取 $L=2.2\Um$，解得 $h=2.1\Um$。（本题为估算题，在取运动员实际长度时可以有一个范围，但要符合实际，故求得 $h$ 值可以不同，$1.6\sim2.6\Um$ 均可）
	\onepicture[w=4.4cm, align=right]{figs/12c_an.jpg}

	\textbf{C．}（1）发生核反应前后，粒子的质量数和核电荷数均不变，据此可知中微子的质量数和电荷数都是 0，A 项正确。

	（2）产生的能量是由于质量亏损。两个电子转变为两个光子之后，质量变为零，由 $E=\Delta mc^{2}$，故一个光子的能量为 $\frac{E}{2}$，代入数据得 $\frac{E}{2}=8.2\times10^{-14}\UJ$。正电子与水中的电子相遇，与电子形成几乎静止的整体，故系统总动量为零，故如果只产生一个光子是不可能的，因为此过程遵循动量守恒。

	（3）物质波的波长为 $\lambda=\frac{h}{p}$，要比较波长需要将中子和电子的动量用动能表示出来即 $p=\sqrt{2mE_{k}}$，因为 $m_{e}<m_{n}$，所以 $p_{e}<p_{n}$，故 $\lambda_{n}<\lambda_{e}$。
}

\item
%% number: 13
%% paperName: 2009年江苏高考第 13 题 15 分
%% typeId: 6
%% type: 计算题
%% chapter: 运动和力
%% point: 加速减速模型
%% method: 无
%% score: 15
%% degree: 450
%% duplicateId: 0
%% body:
航模兴趣小组设计出一架遥控飞行器，其质量 $m=2\Ukg$，动力系统提供的恒定升力 $F=28\UN$。试飞时，飞行器从地面由静止开始竖直上升。设飞行器飞行时所受的阻力大小不变，$g$ 取 $10\Umsq$。
\begin{enumerate}
	\item
	第一次试飞，飞行器飞行 $t_{1}=8\Us$ 时到达高度 $H=64\Um$。求飞行器所受阻力 $f$ 的大小；
	\item
	第二次试飞，飞行器飞行 $t_{2}=6\Us$ 时遥控器出现故障，飞行器立即失去升力。求飞行器能达到的最大高度 $h$；
	\item
	为了使飞行器不致坠落到地面，求飞行器从开始下落到恢复升力的最长时间 $t_{3}$。
\end{enumerate}
%% answer:
\jdanswer{
\begin{enumerate}
	\item
	$f=4\UN$
	\item
	$h=42\Um$
	\item
	$t_{3}=\frac{3\sqrt{2}}{2}\Us$（或 $2.1\Us$）
\end{enumerate}
}
%% memo:
\memoanswer{
	（1）第一次飞行中，设加速度为 $a_{1}$，匀加速运动 $H=\frac{1}{2}a_{1}t_{1}^{2}$，由牛顿第二定律 $F-mg-f=ma_{1}$，解得 $f=4\UN$。

	（2）第二次飞行中，设失去升力时的速度为 $v_{1}$，上升的高度为 $s_{1}$，匀加速运动 $s_{1}=\frac{1}{2}a_{1}t_{2}^{2}$；设失去升力后的加速度为 $a_{2}$，上升的高度为 $s_{2}$，由牛顿第二定律 $mg+f=ma_{2}$，$v_{1}=a_{1}t_{2}$，$s_{2}=\frac{v_{1}^{2}}{2a_{2}}$，解得 $h=s_{1}+s_{2}=42\Um$。

	（3）设失去升力下降阶段加速度为 $a_{3}$；恢复升力后加速度为 $a_{4}$，恢复升力时速度为 $v_{3}$，由牛顿第二定律 $mg-f=ma_{3}$，$F+f-mg=ma_{4}$，且 $\frac{v_{3}^{2}}{2a_{3}}+\frac{v_{3}^{2}}{2a_{4}}=h$，$v_{3}=a_{3}t_{3}$，解得 $t_{3}=\frac{3\sqrt{2}}{2}\Us$（或 $2.1\Us$）。

	运动过程如图所示。
	\onepicture[w=3.4cm, align=right]{figs/13_an.jpg}
}

\item
%% number: 14
%% paperName: 2009年江苏高考第 14 题 16 分
%% typeId: 6
%% type: 计算题
%% chapter: 磁场
%% point: 带电粒子在组合场中的运动
%% method: 无
%% score: 16
%% degree: 450
%% duplicateId: 0
%% body:
1932 年，劳伦斯和利文斯设计出了回旋加速器。回旋加速器的工作原理如图所示，置于高真空中的 D 形金属盒半径为 $R$，两盒间的狭缝很小，带电粒子穿过的时间可以忽略不计。磁感应强度为 $B$ 的匀强磁场与盒面垂直。A 处粒子源产生的粒子，质量为 $m$、电荷量为 $+q$，在加速器中被加速，加速电压为 $U$。加速过程中不考虑相对论效应和重力作用。
\begin{enumerate}
	\item
	求粒子第 2 次和第 1 次经过两 D 形盒间狭缝后轨道半径之比；
	\item
	求粒子从静止开始加速到出口处所需的时间 $t$；
	\item
	实际使用中，磁感应强度和加速电场频率都有最大值的限制。若某一加速器磁感应强度和加速电场频率的最大值分别为 $B_{m}$、$f_{m}$，试讨论粒子能获得的最大动能 $E_{km}$。
\end{enumerate}
\onepicture[w=6.5cm, align=right]{figs/14.png}
%% answer:
\jdanswer{
\begin{enumerate}
	\item
	$r_{2}:r_{1}=\sqrt{2}:1$
	\item
	$t=\frac{\pi BR^{2}}{2U}$
	\item
	$E_{km}=\frac{q^{2}B_{m}^{2}R^{2}}{2m}$ 或 $E_{km}=2\pi^{2}mf_{m}^{2}R^{2}$
\end{enumerate}
}
%% memo:
\memoanswer{
	（1）设粒子第 1 次经过狭缝后的半径为 $r_{1}$，速度为 $v_{1}$，$qU=\frac{1}{2}mv_{1}^{2}$，$qv_{1}B=m\frac{v_{1}^{2}}{r_{1}}$，解得 $r_{1}=\frac{1}{B}\sqrt{\frac{2mU}{q}}$；同理，粒子第 2 次经过狭缝后的半径 $r_{2}=\frac{1}{B}\sqrt{\frac{4mU}{q}}$，则 $r_{2}:r_{1}=\sqrt{2}:1$。

	（2）设粒子到出口处被加速了 $n$ 圈，$2nqU=\frac{1}{2}mv^{2}$，$qvB=m\frac{v^{2}}{R}$，$T=\frac{2\pi m}{qB}$，$t=nT$，解得 $t=\frac{\pi BR^{2}}{2U}$。

	（3）加速电场的频率应等于粒子在磁场中做圆周运动的频率，即 $f=\frac{qB}{2\pi m}$；当磁感应强度为 $B_{m}$ 时，加速电场的频率应为 $f_{Bm}=\frac{qB_{m}}{2\pi m}$，粒子的动能 $E_{K}=\frac{1}{2}mv^{2}$。当 $f_{Bm}\leq f_{m}$ 时，粒子的最大动能由 $B_{m}$ 决定，$qv_{m}B_{m}=m\frac{v_{m}^{2}}{R}$，解得 $E_{km}=\frac{q^{2}B_{m}^{2}R^{2}}{2m}$；当 $f_{Bm}\geq f_{m}$ 时，粒子的最大动能由 $f_{m}$ 决定，$v_{m}=2\pi f_{m}R$，解得 $E_{km}=2\pi^{2}mf_{m}^{2}R^{2}$。
}

\item
%% number: 25
%% paperName: 2009年江苏高考第 25 题 16 分
%% typeId: 6
%% type: 计算题
%% chapter: 电磁感应
%% point: 电磁感应能量分析
%% method: 微元
%% score: 16
%% degree: 250
%% duplicateId: 0
%% body:
如图所示，两平行的光滑金属导轨安装在一光滑绝缘斜面上，导轨间距为 $l$、足够长且电阻忽略不计，导轨平面的倾角为 $\alpha$，条形匀强磁场的宽度为 $d$，磁感应强度大小为 $B$、方向与导轨平面垂直。长度为 $2d$ 的绝缘杆将导体棒和正方形的单匝线框连接在一起组成“$\overline{D}$”型装置，总质量为 $m$，置于导轨上。导体棒中通以大小恒为 $I$ 的电流（由外接恒流源产生，图中未画出）。线框的边长为 $d$（$d<l$），电阻为 $R$，下边与磁场区域上边界重合。将装置由静止释放，导体棒恰好运动到磁场区域下边界处返回，导体棒在整个运动过程中始终与导轨垂直。重力加速度为 $g$。求：
\begin{enumerate}
	\item
	装置从释放到开始返回的过程中，线框中产生的焦耳热 $Q$；
	\item
	线框第一次穿越磁场区域所需的时间 $t_{1}$；
	\item
	经过足够长时间后，线框上边与磁场区域下边界的最大距离 $x_{m}$。
\end{enumerate}
\onepicture[w=7.0cm, align=right]{figs/25.png}
%% answer:
\jdanswer{
\begin{enumerate}
	\item
	$Q=4mgd\sin\alpha-BIld$
	\item
	$t_{1}=\frac{\sqrt{2m(BIld-2mgd\sin\alpha)}+\frac{2B^{2}d^{3}}{mR}}{mg\sin\alpha}$
	\item
	$x_{m}=\frac{BIld}{BIl-mg\sin\alpha}$
\end{enumerate}
}
%% memo:
\memoanswer{
	（1）设装置由静止释放到导体棒运动到磁场下边界的过程中，作用在线框上的安培力做功为 $W$，由动能定理 $mg\sin\alpha\cdot4d+W-BIld=0$，且 $Q=-W$，解得 $Q=4mgd\sin\alpha-BIld$。

	（2）设线框刚离开磁场下边界时的速度为 $v_{1}$，则接着向下运动 $2d$，由动能定理 $mg\sin\alpha\cdot2d-BIld=0-\frac{1}{2}mv_{1}^{2}$。装置在磁场中运动时受到的合力 $F=mg\sin\alpha-F'$，感应电动势 $E=Bdv$，感应电流 $I'=E/R$，安培力 $F'=BI'd$。由牛顿第二定律，在 $t$ 到 $t+\Delta t$ 时间内，有 $\Delta v=\frac{F}{m}\Delta t$，则 $\sum\Delta v=\sum\left[g\sin\alpha-\frac{B^{2}d^{2}v}{mR}\right]\Delta t$，有 $v_{1}=gt_{1}\sin\alpha-\frac{2B^{2}d^{3}}{mR}$，解得 $t_{1}=\frac{\sqrt{2m(BIld-2mgd\sin\alpha)}+\frac{2B^{2}d^{3}}{mR}}{mg\sin\alpha}$。

	（3）经过足够长时间后，线框在磁场下边界与最大距离 $x_{m}$ 之间往复运动，由动能定理 $mg\sin\alpha\cdot x_{m}-BIl(x_{m}-d)=0$，解得 $x_{m}=\frac{BIld}{BIl-mg\sin\alpha}$。
}

\end{enumerate}

\end{document}
